% titel: Blatt 4, A4: de-Broglie-Wellenlänge und Heisenberg-Unschärfe
% skript: Mortimer Kap. 6 (Materiewellen, Unschärferelation)

% Teil I: Werkzeuge, kein Lösungsweg.

Zwei grundlegende Wellenphänomene der Materie: Jede Masse hat eine
de-Broglie-Wellenlänge, und Ort und Impuls sind prinzipiell nicht gleichzeitig
beliebig genau bestimmbar.

\begin{werkzeugkasten}[de-Broglie-Wellenlänge]
Jedes Teilchen mit Impuls $p = m v$ hat die Wellenlänge
\begin{equation}
  \lambda = \frac{h}{p} = \frac{h}{m\,v}.
\end{equation}
Für makroskopische Körper ist $\lambda$ verschwindend klein (kein
beobachtbarer Wellencharakter), für Elektronen atomare Größen.

\textbf{Heisenberg-Unschärfe:}
\begin{equation}
  \Delta x \cdot \Delta p \ge \frac{\hbar}{2},
  \qquad \Delta p = m\,\Delta v
  \;\Longrightarrow\;
  \Delta v \ge \frac{\hbar}{2\,m\,\Delta x}.
\end{equation}
$\hbar = h/(2\pi) \approx \qty{1.055e-34}{J\,s}$.

\textit{Skript:} Mortimer Kap. 6 (Materiewellen, Unschärferelation),
Übung 4/4.
\end{werkzeugkasten}

\begin{achtung}[Darauf achten]
\begin{itemize}
  \item Geschwindigkeit in m/s umrechnen (130 km/h $=\qty{36.11}{m/s}$,
        $\;c/1000 \approx \qty{2.998e5}{m/s}$).
  \item Masse in kg (1.35 t $= \qty{1350}{kg}$; $m_e = \qty{9.109e-31}{kg}$).
  \item $\Delta x$ in m (10 pm $= 10^{-11}\ \mathrm{m}$); mit $\hbar=h/2\pi$
        rechnen (Faktor $4\pi$ im Nenner, wenn mit $h$ statt $\hbar$
        gearbeitet wird).
  \item Interpretation: kleine Masse $\to$ grosse $\lambda$ und grosse
        Mindest-Unschärfe in $\Delta v$.
\end{itemize}
\end{achtung}
